% =====================================================================
% Secao de Evaluation reescrita (metodologia + resultados experimentais).
% Substitui/expande a atual "Evaluation Methodology" do main.tex.
% Idioma: ingles, para casar com o restante do artigo.
% Insira as tabelas de results_tables.tex onde indicado.
% =====================================================================

\section{Evaluation}
\label{sec:evaluation}

The architectural argument of this paper is comparative: the compartmentalized
design should contain compromise better than a centralized baseline while paying
an acceptable operational cost. We therefore evaluate the implementation against
four explicit criteria derived from the design goals: (i)~\emph{compromise
containment}, whether a failure in one plane reaches assets owned by another;
(ii)~\emph{exfiltration resistance}, whether execution-time code can open
outbound channels; (iii)~\emph{persistence resistance}, whether one run can
contaminate later runs; and (iv)~\emph{operational cost}, the latency introduced
by the isolation strategy.

\subsection{Methodology}

We exercise the running prototype with a small suite of purpose-built
applications, each designed to attack one criterion. The applications and the
synthetic, PII-free dataset used as input are versioned with the artifact
repository (\texttt{assets/apps/} and \texttt{assets/data/}). Tests that target
the execution plane submit an application bundle directly to the executor service
(\texttt{POST /execute}), reproducing exactly the \texttt{workspace/run.sh}
contract that the control plane generates; this isolates the variable under study
(the execution microVM) from unrelated services. Tests that target lateral
movement inspect the executor container, and the access-control test drives the
public marketplace API. Every adverse application reports machine-readable
evidence markers on standard output, which the harness
(\texttt{assets/tests/}) collects into timestamped evidence files. Each test is
mapped to the design claim it is meant to support, following the principle that
an evaluation should connect each observation to a property rather than merely
listing successful checks.

The centralized baseline (Environment~A) keeps execution under the same
Firecracker isolation as the compartmentalized design (Environment~B), so the
only independent variable between the two is the separation of planes. This makes
the comparison fair: any measured difference reflects compartmentalization, not a
weaker execution sandbox.

\subsection{Containment, Exfiltration, and Persistence}

Table~\ref{tab:security-results} summarizes the suite. The execution plane shows
the strongest, most direct evidence. For exfiltration resistance (T1), the
malicious application observes only the loopback and \texttt{sit0} pseudo-devices
inside the guest: there is no \texttt{eth0}, because the execution microVM is
launched without a network interface. All seven outbound attempts---DNS
resolution, raw UDP and TCP to public resolvers, HTTPS, plain HTTP, and an SMTP
probe---fail with \texttt{Network is unreachable} or name-resolution errors. The
channel is closed by construction rather than by a firewall rule that an attacker
might disable, which is the strongest form the property can take.

For containment at the filesystem level (T1b), the execution code runs as root
inside the disposable guest but finds no control-plane secrets in its environment
(\texttt{ASSET\_ENCRYPTION\_KEY}, \texttt{PRIVATE\_KEY}, and related variables are
all absent), and it can enumerate only its own input and output disks; sibling
execution disks are not present. The files it can read
(\texttt{/etc/passwd}, \texttt{/etc/shadow}) belong to the minimal ephemeral guest
rootfs, not to the host, and contain no credential material. Writes outside the
workspace succeed only on ephemeral overlay disks that are destroyed at teardown.

That teardown is confirmed by the persistence test (T3): after more than nine
consecutive executions, the executor retains no Firecracker socket, no
\texttt{rootfs-overlay}, \texttt{input}, \texttt{output}, or \texttt{application}
\texttt{ext4} image, no residual loopback mount, and no leftover dataset or
working directory. The \texttt{cleanup} trap removes the entire per-run working
directory, so no ephemeral state survives to contaminate a later run.

The supply-chain test (T2) connects the two planes. We seeded a malicious package
into the builder's wheel cache and registered an application that declares it as a
dependency. The builder resolved and bundled the package (a 137~s build that
reported the cache hit), and the package then ran its payload at import time
inside the execution microVM. That payload is the interesting evidence: the
dependency's code does execute, but it finds no control-plane secrets in the
environment and reports its outbound probe to a public resolver as blocked,
exactly mirroring the T1 and T1b results. A compromised dependency therefore
inherits the containment of the plane that runs it. The complementary build-time
vector (a source distribution whose \texttt{setup.py} runs during the build)
begins executing inside the build microVM before failing under the read-only
cache mount, which confirms that such code, when it does run, runs in the
isolated builder rather than in the control plane.

\subsection{Blast Radius and Lateral Movement}

The clearest architectural difference between the two designs appears in the
cross-plane reachability test (T4), shown in Table~\ref{tab:blast-radius}. From a
hypothetically compromised executor in Environment~B, the asset-encryption key
and the wallet private key are absent from the environment, and the IPFS
repository, the wheel cache, and the executions store are absent from the
filesystem. The same assets are, by construction of the centralized image,
colocated in the single namespace of Environment~A: there the very secrets and
stores that are unreachable in B sit beside the execution logic.

We report one residual weakness honestly. In the single-host local deployment,
all three roles share one Docker network, so the executor can still reach the
control API (port~3001) and the blockchain RPC (port~8545) over that network,
even though the IPFS API (port~5001), which would expose the encrypted blobs
directly, is bound to loopback inside the control plane and is therefore
unreachable. This substantiates---rather than contradicts---the limitation
already stated in the threat model: internal service authentication and network
segmentation (mutual TLS) remain future work. In a distributed deployment across
separate virtual machines, reachability of ports~3001 and 8545 becomes a matter
of inter-host network policy. The key point for the comparative argument stands:
the secrets and persistent stores that define the blast radius are isolated in B
and colocated in A.

\subsection{Access Control}

The access-control test (T5) confirms that execution is gated by the smart
contract. After registering the synthetic dataset on chain---whose on-chain
content hash matches the file's SHA-256, evidence of integrity binding---an
execution request issued without a prior \texttt{purchaseAccess} call is rejected
with \emph{access denied} by the \texttt{hasAccessByIds} check, before any
execution bundle is prepared. The control plane therefore refuses to move data or
code toward the executor for an unauthorized caller.

\subsection{Operational Cost}

Table~\ref{tab:performance} reports the cost of the isolation strategy on the
execution path. Across six consecutive runs of the benign application over the
synthetic dataset, the end-to-end execution latency averages 31.4~s
(median 31.8~s, standard deviation 1.4~s, range 28.5--32.9~s). This time is
dominated by the microVM life cycle---kernel boot, disk discovery, output-disk
synchronization, and the guarded shutdown sequence---rather than by the
application itself, whose computation over sixty rows completes in milliseconds.
The roughly thirty-second fixed overhead per execution is the price of launching
a fresh, networkless microVM for every run; it is the operational counterpart of
the containment and persistence properties demonstrated above.

The build plane is more expensive, as expected from its heavier configuration
(two vCPUs, more memory, network egress, and dependency resolution). A
dependency-free build completes in about 103~s and produces a 16~MiB artifact; a
build that pulls a real dependency tree (\texttt{requests} and four transitive
packages) takes about 131~s from a cold cache and grows the artifact to 19.6~MiB.
The wheel cache is a measurable optimization: a warm-cache rebuild of the same
dependency set drops to about 114~s and avoids re-downloading, confirming that
caching trades a controlled supply-chain surface for build-time savings. Build
latencies are single-run samples with higher variance than the execution
measurements, but the order of magnitude is stable: building is roughly three to
four times the cost of an execution.

\subsection{Threats to Validity}

The measurements were taken on a single host with a Linux-backed container
runtime; absolute latencies will vary with hardware, but the containment and
exfiltration results are structural (absence of a network interface, absence of
secrets) and do not depend on timing. The execution-plane tests bypass IPFS and
the builder to isolate the microVM under study; the end-to-end path through the
control plane is exercised separately by the access-control test. Finally, the
baseline comparison for blast radius is derived from the deployment configuration
rather than from a live intrusion, because both environments share the same
execution sandbox and differ only in colocation of assets.
